\documentclass[10pt,a4paper]{amsart}
\usepackage[margin=19mm]{geometry}
\usepackage{amsmath,amssymb,mathtools}
\usepackage[scr=rsfs,cal=euler]{mathalfa}
\usepackage{xfrac}
\usepackage[unicode,hidelinks,bookmarks=false]{hyperref}
\newcommand{\ie}{i.e.\ }
\newcommand{\cf}{cf.\ }
\newcommand{\resp}{respectively\ }
\newcommand{\Subsection}{\subsection}
\providecommand{\nobreakdash}{\nobreak\hbox{-}}
\setlength{\emergencystretch}{2em}
\multlinegap0pt
\usepackage{bm}
\usepackage{bbm}
\usepackage{dsfont}
\usepackage{xspace}
\usepackage{cancel}
\usepackage{mathtools}
\usepackage{amsthm}
\makeatletter
\@ifundefined{theorem}{\newtheorem{theorem}{Theorem}}{}
\makeatother
\title[On a generalization of Roman domination with more legions]{On a generalization of Roman domination with more legions}
\author{Fahimeh Khosh-Ahang Ghasr}
\date{Source: arXiv:2212.08935v1 (CC BY 4.0)}
\begin{document}
\maketitle

\noindent\emph{Source and licence.} On a generalization of Roman domination with more legions. Version arXiv:2212.08935v1 (CC BY 4.0), distributed under the Creative Commons Attribution 4.0 International licence, as recorded in the arXiv licence metadata for this exact version and shown on the abstract page https://arxiv.org/abs/2212.08935v1. Full rights notice: licenses/arxiv-2212.08935-CC-BY-4.0.txt.\par\smallskip

\noindent\textit{Editorial scope.} Counted unit: the Theorem whose source label is \texttt{thm4} in the article (source0.tex lines 773--788, source label \texttt{thm4}), delivered together with the article's own complete original proof of it (source0.tex lines 789--958, one \texttt{proof} environment of 170 lines). No mathematics is written, repaired or re-derived: statement and proof are the article's own text line for line. Required context retained uncounted: none. Every definition, display and label of those lines is retained unchanged. Omitted from this package and not redistributed: 1--772, 959--1016. Nothing inside a retained excerpt is omitted or altered: every retained body is delivered exactly as the article's lines are written, so no omission receipt of any kind accompanies this package. Citations: the retained excerpts carry no citation command at all, so no bibliography entry of the article is transcribed and no reference is redistributed. Reference adaptation: none was needed and none was made; every reference inside the delivered unit resolves inside it, so no unresolved reference or citation is produced. Printed slips kept literal: no printed word, symbol, hypothesis or number of the counted statement or of its proof was altered. Layout and class: the article's own class is \texttt{amsart} with packages including amscd, amsmath, amsthm, amssymb, lineno, pstcol, pst-plot, pst-3d, color, pstricks, stmaryrd, inputenc, pstricks-add, graphicx; the wrapper is the lane's pinned \texttt{amsart} editorial wrapper at 10pt on A4, which restates the theorem-like environments the retained text needs and adds the article's own macro definitions that the retained text needs (none), with extra wrapper packages (bm, bbm, dsfont, xspace, cancel, mathtools). The delivered numbering is the unit's own. Assets: no retained span contains a figure environment, a graphics-inclusion command, a PDF-page inclusion, a table or a TikZ picture; no image file is copied into this package, the package declares no asset dependency and no image file is redistributed with it. Licence: version arXiv:2212.08935v1 of this article is distributed under the Creative Commons Attribution 4.0 International licence, as recorded in the arXiv licence metadata for this exact version and shown on the abstract page https://arxiv.org/abs/2212.08935v1. A full rights notice is delivered at licenses/arxiv-2212.08935-CC-BY-4.0.txt. Characters: every retained excerpt is pure ASCII, so the delivered engine, XeLaTeX with its default Latin Modern text font, reports no missing character.
\begin{theorem}\label{thm4}
$$\gamma_k(K_{m,n})=\left\lbrace \begin{array}{ll}
k; & m=1, n\geq 1 \\ 
k; & k \textrm{ is even}, m=2, n\geq 2 \\
k+1; &  k \textrm{ is odd, } m=2, n\geq 2 \\
\frac{3k}{2}; & k \textrm{ is even}, m=3, n\geq 3 \\
\frac{3k+1}{2}; & k=1,3 , m=3, n\geq 3 \\  
\frac{3k+1}{2}; & k \textrm{ is odd}, m=n=3 \\
\frac{3k+3}{2}; & 5\leq k \textrm{ is odd}, m=3, n\geq 6 \\
\frac{3k+3}{2}; & 11\leq k \textrm{ is odd}, m=3, n=5 \\
\frac{3k+3}{2}; & 17\leq k \textrm{ is odd}, m=3, n=4 \\
\frac{3k+1}{2} \textrm{ or } \frac{3k+3}{2}; & 5\leq k \leq 15, k \textrm{ is odd}, m=3, n=4 \\
\frac{3k+1}{2} \textrm{ or } \frac{3k+3}{2}; & 5\leq k \leq 9, k \textrm{ is odd}, m=3, n=5 \\
\frac{2m}{m+1}k \leq \gamma_k(K_{m,n})\leq 2k; & n\geq m\geq 4 
\end{array} \right.$$
\end{theorem}

\begin{proof}
We separate the proof in different situations as follows:

\textbf{A.  ${\bf m=1}$ and ${\bf n\geq 1}$.} 

By assigning $k$ to the vertex of $A$ and zero to other vertices we get to a Roman $k$-dominating function with weight $k$. So $\gamma_k (K_{1,n})=k$. 

\textbf{B. ${\bf k}$ is even, ${\bf m=2}$ and ${\bf n\geq 2}$.} 

By assigning $k/2$ to the vertices of $A$ and zero to other vertices we get to a Roman $k$-dominating function with weight $k$. So $\gamma_k (K_{2,n})=k$ for even integers $k$. 

Hereafter suppose that $f$ is a Roman $k$-dominating function.

\textbf{C. ${\bf k}$ is odd, ${\bf m=2}$ and ${\bf n\geq 2}$.} 

Case I. $f(u)\geq (k+1)/2$ for all $u\in A$ (or similarly $f(v)\geq (k+1)/2$ for all $v\in B$). Then since $|A|\geq 2$, clearly we have $w(f)\geq k+1$.

Case II.  $f(u_1)<k/2$ and $f(u_2)\geq (k+1)/2$ (or similarly $f(v_1)<k/2$ and $f(v_2)\geq (k+1)/2$ for some $v_1, v_2\in B$). Then 
\begin{align*}
w(f)&\geq f(u_2)+(f(u_1)+\sum_{v\in B}f(v))\\
&\geq (k+1)/2+k\\
&\geq k+1.
\end{align*}

Case III.  $f(u)<k/2$ for all $u\in V$. Therefore
$$f(u_1)+\sum_{v\in B}f(v)\geq k, f(u_2)+\sum_{v\in B}f(v)\geq k,$$
$$f(v_1)+\sum_{u\in A}f(u)\geq k, f(v_2)+\sum_{u\in A}f(u)\geq k,$$
for some $v_1, v_2\in B$.
Summing up the above inequalities implies that $3w(f)\geq 4k$, which ensures that $w(f)\geq k+1$, since $w(f)$ is an integer.

The above cases show that $\gamma_k(K_{2,n})\geq k+1$. By assigning $(k+1)/2$ to the vertices of $A$ and zero to the vertices of $B$, we get to a Roman $k$-dominating function of $K_{m,n}$ with weight $k+1$. Hence $\gamma_k(K_{2,n})=k+1$,
in this case.


\textbf{D. ${\bf n\geq m\geq 3}$.}

Case I.  $f(u)\geq k/2$ for all $u\in A$ (or similarly $f(v)\geq k/2$ for all $v\in B$). Then 
$$w(f)\geq mk/2\geq \frac{2m}{m+1}k.$$
(of course if $k$ is odd, we even have  $w(f)\geq m(k+1)/2$.)

Case II.  $f(u_1)<k/2$ and $f(u_2), f(u_3)\geq k/2$ for some $u_1, u_2, u_3\in A$ (or similarly  $f(v_1)<k/2$ and  $f(v_2), f(v_3)\geq k/2$ for some $v_1, v_2,v_3\in B$). Then 
$$w(f)\geq (f(u_1)+\sum_{v\in B}f(v))+f(u_2)+f(u_3)\geq k+k/2+k/2=2k.$$

Case III. $f(u_1)\geq k/2$ and $f(u)<k/2$ for all $u\in V\setminus \{u_1\}$ (or similarly $f(v_1)\geq k/2$ and $f(u)<k/2$ for all $u\in V\setminus \{v_1\}$). Then 
$$f(u_i)+\sum_{v\in B}f(v)\geq k, \ \forall i \textrm{ with }   2\leq i \leq m,$$
$$f(v_i)+\sum_{u\in A}f(u)\geq k, \ \forall i \textrm{ with }  1\leq i \leq m.$$ 
Note that since there exists $u_i\in A$ with $f(u_i)<k/2$ and $f(u_i)+\sum_{v\in B}f(v)\geq k$, we should have $\sum_{v\in B}f(v)\geq k/2$. So, summing up the above inequalities implies $(m+1)w(f)\geq (2m-1)k+f(u_1)+\sum_{v\in B}f(v)\geq 2mk$. Hence
$$w(f)\geq \frac{2m}{m+1}k.$$

Case IV. $f(u_1)\geq k/2, f(v_1)\geq k/2$ for some $u_1\in A, v_1\in B$ and $f(u)<k/2$ for all $u\in V\setminus \{u_1, v_1\}$. Then
$$f(u_i)+\sum_{v\in B}f(v)\geq k, \ \forall i \textrm{ with }  2\leq i \leq m,$$
$$f(v_i)+\sum_{u\in A}f(u)\geq k, \ \forall i \textrm{ with }  2\leq i \leq m.$$
Summing up the above inequalities implies 
$$mw(f)\geq (2m-2)k+f(u_1)+f(v_1)\geq (2m-1)k.$$
So
$$w(f)\geq \frac{2m-1}{m}k> \frac{2m}{m+1}k.$$

Case V. $f(u)<k/2$ for all $u\in V$. Then 
$$f(u_i)+\sum_{v\in B}f(v)\geq k, \ \forall i \textrm{ with }  1\leq i \leq m,$$
$$f(v_i)+\sum_{u\in A}f(u)\geq k, \ \forall i \textrm{ with }  1\leq i \leq m.$$
Summing up the above inequalities implies $(m+1)w(f)\geq 2mk$. So
$$w(f)\geq \frac{2m}{m+1}k.$$

The above cases show that $\gamma_k(K_{m,n})\geq \frac{2m}{m+1}k$ and so 
\begin{equation}
\frac{2m}{m+1}k \leq \gamma_k(K_{m,n})\leq 2k.
\end{equation}

Note that for special values of $k$, we may have $w(f)=\frac{2m}{m+1}k$. For instance, when $k=(m+1)r$ for some integer $r$, by assigning $r$ to each vertex of $K_{m,m}$ we get to a Roman $k$-dominating function with weight $\frac{2m}{m+1}k$.

If $m=3$, then $(5)$ shows that $w(f)\geq 3k/2$. Also if $k$ is even, then by assigning $k/2$ to the vertices of $A$ and zero to the vertices of $B$, one can get to a  Roman $k$-dominating function of $K_{m,n}$ with weight $3k/2$. So 
$$\gamma_k(K_{3,n})=3k/2,$$
when $n\geq 3$ and $k$ is even.

If $m=3$ and $k=1$, then since $w(f)$ is an integer, $(5)$ implies that 
$$\gamma_k(K_{m,n})=(3k+1)/2=2k.$$

If $m=3$ and $1<k$ is odd, then by assigning $(k+1)/2$ to the vertices of $A$ and zero to the vertices of $B$, we have a Roman $k$-dominating function with weight $3(k+1)/2$. Thus $\gamma_k(K_{m,n}) \leq (3k+3)/2 $. So, by $(5)$  for $m=3$, we have 
$$3k/2 \leq \gamma_k(K_{m,n})\leq (3k+3)/2.$$
 Hence since $w(f)$ is an integer, we have $\gamma_k(K_{m,n})=(3k+1)/2$ or $\gamma_k(K_{m,n})=(3k+3)/2$. 

When $m=n=3$, by setting
 $$f(u_1)=(k-3)/4, f(u_2)=f(u_3)=f(v_1)=f(v_2)=f(v_3)=(k+1)/4,$$
 when $k\overset{4}{\equiv} 3$  and 
$$f(u_1)=f(v_1)=(k+3)/4, f(u_2)=f(u_3)=f(v_2)=f(v_3)=(k-1)/4,$$
 when $k\overset{4}{\equiv}1$,  we get to a Roman $k$-dominating function with weight $(3k+1)/2$. Thus 
$$\gamma_k(K_{3,3})=(3k+1)/2.$$
Also, when $k=3$ and $n\geq m=3$, if $f(u)=1$ for all $u\in A$, $f(v_1)=2$ for some $v_1\in B$ and $f(v)=0$ for all $v\in B\setminus \{v_1\}$, then $f$ is a Roman $3$-dominating function with weight $5$ and so 
$$\gamma_3(K_{3,n})=(3k+1)/2.$$
Now suppose that $k>3$ and $n>m=3$. Also assume that there exists a Roman $k$-dominating function, say $f$, with $w(f)=(3k+1)/2$. As the above cases illustrate, it may exists in Cases III or V. 

Firstly in Case III, suppose that $f(v_n)\geq (k+1)/2$ and $f(u)<k/2$ for all $u\in V\setminus \{v_n\}$. Then
$$f(u_i)+\sum_{v\in B}f(v)\geq k, \ \forall i \textrm{ with }   1\leq i \leq 3, $$
$$f(v_i)+\sum_{u\in A}f(u)\geq k, \ \forall i \textrm{ with }   1\leq i \leq n-1.$$
Summing the above inequalities implies that
$$3w(f)+(w(f)-f(v_n))+(n-4)\sum_{u\in A}f(u)\geq (n+2)k.$$
Since $f(u_1)\leq (k-1)/2$, $f(u_1)+\sum_{v\in B}f(v)\geq k$ ensures that $\sum_{v\in B}f(v)\geq (k+1)/2$. Hence
\begin{align*}
nw(f)&\geq (n+2)k+f(v_n)+(n-4)\sum_{v\in B}f(v)\\
&\geq (n+2)k+(k+1)/2+(n-4)(k+1)/2\\
&=(3nk+n+k-3)/2.
\end{align*}
Therefore $k>3$ implies
$$\frac{3k+1}{2}=w(f)\geq \frac{3k+1}{2}+\frac{k-3}{2n}> \frac{3k+1}{2},$$
which is a contradiction.

Now suppose that in Case III, $f(u_3)\geq (k+1)/2$ and $f(u)<k/2$ for all $u\in V\setminus \{v_n\}$. Then
$$f(u_i)+\sum_{v\in B}f(v)\geq k, \ \forall i \textrm{ with }   1\leq i \leq 2, $$
$$f(v_j)+\sum_{u\in A}f(u)\geq k, \ \forall j \textrm{ with }   1\leq j \leq n.$$
Summing the above inequalities for each $1\leq i \leq 2$ and $1\leq j \leq n$ implies that $f(u_i)+f(v_j)+w(f)\geq 2k$. So, $w(f)=(3k+1)/2$ yields to
\begin{equation}\label{k-1}
f(u_i)+f(v_j)\geq 2k-(3k+1)/2=(k-1)/2. 
\end{equation}
Hence if $y=\min\{f(v) \ | \ v\in B\}$, then
\begin{align*}
(3k+1)/2=w(f)&=(f(u_1)+f(v_1))+(f(u_2)+f(v_2))+f(u_3)+\sum_{i=3}^nf(v_i)\\
&\geq (k-1)/2+(k-1)/2+(k+1)/2+(n-2)y\\
&=(3k-1)/2+(n-2)y.
\end{align*}
Thus $y\leq 1/(n-2)$, which implies that $y=0$, since $n>3$. Now, (\ref{k-1}) and $f(u_1), f(u_2)\leq (k-1)/2$ yield to $f(u_1)=f(u_2)=(k-1)/2$. So since $f(u_1)+\sum_{v\in B}f(v)\geq k$, we should have $\sum_{v\in B}f(v)\geq (k+1)/2$. Therefore
\begin{align*}
(3k+1)/2=w(f)&=f(u_1)+f(u_2)+f(u_3)+\sum_{v\in B}f(v)\\
&\geq (k-1)/2+(k-1)/2+(k+1)/2+(k+1)/2\\
&=2k,
\end{align*}
which is a contradiction. So, such $f$ can't be gained by Case III.

Moreover, in Case V we have $f(u)<k/2$ for all $u\in V$. Then
$$f(u_i)+\sum_{v\in B}f(v)\geq k, \ \forall i \textrm{ with }  1\leq i \leq 3,$$
$$f(v_j)+\sum_{u\in A}f(u)\geq k,  \ \forall j \textrm{ with }  1\leq j \leq n.$$
Summing the above inequalities for each $1\leq i \leq 3$ and $1\leq j \leq n$ implies that $f(u_i)+f(v_j)+w(f)\geq 2k$. So, $w(f)=(3k+1)/2$ yields to
\begin{equation}\label{k-2}
f(u_i)+f(v_j)\geq 2k-(3k+1)/2=(k-1)/2. 
\end{equation}
Hence if $y=\min\{f(v) \ | \ v\in B\}$, then
\begin{align*}
(3k+1)/2=w(f)&=(f(u_1)+f(v_1))+(f(u_2)+f(v_2))+(f(u_3)+f(v_3))+\sum_{i=4}^nf(v_i)\\
&\geq 3(k-1)/2+(n-3)y.
\end{align*}
Thus $y\leq 2/(n-3)$. Now if $y=0$, then (\ref{k-2}) and $f(u_i)\leq (k-1)/2$ yield to $f(u_i)=(k-1)/2$ for all $i=1,2,3$. Hence
$$(3k+1)/2=w(f)=\sum_{u\in A}f(u)+\sum_{v\in B}f(v)\geq 3(k-1)/2+(k+1)/2=(4k-2)/2.$$
This implies that $k\leq 3$ which is a contradiction. Hence when $5\leq k$ is an odd number and $n\geq 6$, such $f$ can't be gained by Case V. So
$$\gamma_k(K_{3,n})=(3k+3)/2,$$
when $k\geq 5$ is an odd number, $n\geq 6$ and $m=3$.

Now if $y=1$, then (\ref{k-2}) and $f(u_i)\leq (k-1)/2$ yield to $f(u_i)\geq (k-3)/2$ for all $i=1,2,3$. Hence
$$(3k+1)/2=w(f)=\sum_{u\in A}f(u)+\sum_{v\in B}f(v)\geq 3(k-3)/2+(k+1)/2=(4k-8)/2.$$
This implies that $k\leq 9$. Hence 
$$\gamma_k(K_{3,5})=(3k+3)/2,$$
when $k\geq 11$ is an odd number.

Also, if $y=2$, then similarly we have 
$$(3k+1)/2=w(f)=\sum_{u\in A}f(u)+\sum_{v\in B}f(v)\geq 3(k-5)/2+(k+1)/2=(4k-14)/2.$$
This implies that $k\leq 15$. Hence 
$$\gamma_k(K_{3,4})=(3k+3)/2,$$
when $k\geq 17$ is an odd number.
 
 These complete the proof. Note that as discussed earlier, when either $k$ is odd with  $5\leq k \leq 15$ and $m=3, n=4$ or $k$ is odd with $5\leq k \leq 9$ and $m=3, n=5$, the precise value of $\gamma_k(K_{m,n})$ can be found by Case V and $y=\min\{f(v) \ | \ v\in B\}=1$ or $2$. For instance when $k=5, n=4, m=3$, if $f(u_1)=2$ and $f(v)=1$ for all $v\in V\setminus \{u_1\}$, then  $f$ is a Roman $5$-dominating function with weight $8$ and so 
$$\gamma_5(K_{3,4})=(3k+1)/2.$$
When $k=7, n=5, m=3$, if $f(u)=2$ and $f(v)=1$ for all $u\in A$ and $v\in B$, then  $f$ is a Roman $7$-dominating function with weight $11$ and so 
$$\gamma_7(K_{3,5})=(3k+1)/2.$$
When $k=11, n=4, m=3$, if $f(u)=3$ and $f(v)=2$ for all $u\in A$ and $v\in B$, then  $f$ is a Roman $11$-dominating function with weight $17$ and so 
$$\gamma_{11}(K_{3,4})=(3k+1)/2.$$
But when $k=5, n=5, m=3$, the above discussion shows that $y=1$ and so $f(u_i)=1$ or $2$. Since $y+\sum_{u\in A}f(u)\geq 5$ and $f(v)\geq y$ for all $v\in B$, we have
$$w(f)=\sum_{u\in A}f(u)+\sum_{v\in B}f(v)\geq 9>(3k+1)/2.$$
Hence
$$\gamma_5(K_{3,5})=(3k+3)/2.$$
Therefore, in each case it should be discussed.

\end{proof}

\end{document}
